\chapter{จุลินทรีย์บริเวณรากพืชและการส่งเสริมการเจริญเติบโต}
\label{chap:ch03}

% ==========================================================
% แผนบริหารการสอนประจำบทที่ 3
% ==========================================================
\section*{แผนบริหารการสอนประจำบทที่ 3}
\addcontentsline{toc}{section}{แผนบริหารการสอนประจำบทที่ 3}

\textbf{หัวข้อเนื้อหาประจำบท}
\begin{enumerate}
    \item นิเวศวิทยาบริเวณไรโซสเฟียร์ (Rhizosphere Biology) และผลกระทบจากสารหลั่งรากพืช
    \item แบคทีเรียส่งเสริมการเจริญเติบโตของพืช (PGPR) และการจำแนกประเภท
    \item กลไกทางตรงและทางอ้อมในการส่งเสริมการเติบโตและการทนทานต่อความเครียด
    \item ชีววิทยาของเชื้อราอาร์บัสคูลาร์ไมคอร์ไรซา (AMF)
    \item การพัฒนาและใช้ประโยชน์จากจุลินทรีย์เอนโดไฟต์ (Endophytic Microorganisms)
\end{enumerate}

\textbf{วัตถุประสงค์เชิงพฤติกรรม}
\begin{itemize}
    \item อธิบายความหมายและองค์ประกอบของบริเวณไรโซสเฟียร์ ไรโซเพลน และเอนโดสเฟียร์ได้
    \item วิเคราะห์กลไกการสังเคราะห์ฮอร์โมน IAA และเอนไซม์ ACC deaminase ของ PGPR ได้
    \item อธิบายโครงสร้างและการแลกเปลี่ยนสารอาหารในระบบชีวภาพไมคอร์ไรซาได้
    \item ออกแบบสูตรตำรับหัวเชื้อจุลินทรีย์เพื่อส่งเสริมการแตกรากของกล้าไม้ผลได้
\end{itemize}

\textbf{กิจกรรมการเรียนการสอน}
\begin{itemize}
    \item การบรรยายและชมภาพจำลอง 3 มิติการเข้าอยู่อาศัยของแบคทีเรียบริเวณผิวราก
    \item ปฏิบัติการย้อมสีรากพืชเพื่อตรวจหาโครงสร้าง Arbuscules ของเชื้อรา AMF
    \item การตรวจวัดการผลิตฮอร์โมนออกซิน (IAA) ของแบคทีเรียด้วยน้ำยา Salkowski Reagent
\end{itemize}

\textbf{สื่อการเรียนการสอน}
\begin{itemize}
    \item เอกสารคำสอนบทที่ 3 และภาพตัดขวางแสดงโครงสร้างทางกายภาพบริเวณรากพืช
    \item ตัวอย่างรากข้าวโพดและรากพืชตระกูลถั่วที่ติดเชื้อไมคอร์ไรซา
\end{itemize}

\textbf{การประเมินผล}
\begin{itemize}
    \item การประเมินทักษะการเตรียมสไลด์รากพืชและบันทึกภาพใต้กล้องจุลทรรศน์
    \item การตอบคำถามเชิงวิเคราะห์กลไกความเครียดและการยืดอายุพืช
\end{itemize}

\newpage

% ==========================================================
% เนื้อหาบทเรียน
% ==========================================================

\section{นิเวศวิทยาบริเวณไรโซสเฟียร์}
\index{ไรโซสเฟียร์ (Rhizosphere)}

คำว่า \textbf{“ไรโซสเฟียร์ (Rhizosphere)”} ได้รับการบัญญัติขึ้นเป็นครั้งแรกในคริสต์ศักราช 1904 โดย \textbf{โลเรนซ์ ฮิลต์เนอร์ (Lorenz Hiltner, 1862--1923)} นักจุลชีววิทยาและพยาธิวิทยาพืชชาวเยอรมัน เพื่ออธิบายถึงบริเวณชั้นดินบางๆ ที่อยู่ติดกับผิวรากพืช (โดยทั่วไปอยู่ในระยะ 1--5 มิลลิเมตรจากผิวราก) ซึ่งได้รับอิทธิพลโดยตรงจากกิจกรรมและสารที่รากพืชปลดปล่อยออกมา

บริเวณรอบรากพืชสามารถแบ่งออกเป็น 3 โซนตามระยะห่างทางกายภาพ:
\begin{enumerate}
    \item \textbf{เอนโดสเฟียร์ (Endosphere)} เนื้อเยื่อภายในรากพืชซึ่งเป็นแหล่งอาศัยของจุลินทรีย์เอนโดไฟต์
    \item \textbf{ไรโซเพลน (Rhizoplane)} ผิวหน้าภายนอกของเซลล์ผิวราก (Epidermis) และขนราก (Root hairs)
    \item \textbf{ไรโซสเฟียร์ (Rhizosphere)} ชั้นดินรอบผิวรากที่อุดมไปด้วยสารหลั่งจากราก
\end{enumerate}

รากพืชจะปลดปล่อยสารอินทรีย์ที่สังเคราะห์ได้จากกระบวนการสังเคราะห์ด้วยแสงออกมาสู่ไรโซสเฟียร์ในรูปของ \textbf{สารหลั่งจากราก (Root Exudates)} คิดเป็นสัดส่วนสูงถึง 10--40\% ของคาร์บอนทั้งหมดที่พืชตรึงได้ สารเหล่านี้ประกอบด้วยน้ำตาล กรดอะมิโน กรดอินทรีย์ วิตามิน และสารทุติยภูมิกลุ่มฟลาโวนอยด์ (Flavonoids) ส่งผลให้ความหนาแน่นของประชากรจุลินทรีย์บริเวณไรโซสเฟียร์สูงกว่าดินทั่วไป (Bulk Soil) มากถึง 10 ถึง 100 เท่า ซึ่งเรียกว่าปรากฏการณ์ \textbf{The Rhizosphere Effect (R:S ratio)}

\section{แบคทีเรียส่งเสริมการเจริญเติบโตของพืช (PGPR)}
\index{แบคทีเรียพีจีพีอาร์ (PGPR)}

\textbf{Plant Growth-Promoting Rhizobacteria (PGPR)} คือกลุ่มแบคทีเรียในดินที่มีความสามารถในการเข้าครอบครองพื้นที่บริเวณรากพืช (Root Colonization) และกระตุ้นการเจริญเติบโตหรือเพิ่มความทนทานต่อสภาวะเครียดของพืช สกุลที่สำคัญได้แก่ \textit{Pseudomonas}, \textit{Bacillus}, \textit{Azospirillum}, \textit{Burkholderia}, \textit{Enterobacter}, และ \textit{Serratia}

\begin{figure}[H]
\centering
\begin{tikzpicture}[scale=0.95]
    % Central root tissue cross-section
    \fill[agriEarth!20, draw=agriEarth!80, line width=1.5pt] (0,-3.5) rectangle (2.5, 3.5);
    \node at (1.25, 0) [rotate=90, font=\footnotesize\bfseries\color{agriEarth!90!black}, align=center] {เนื้อเยื่อรากพืช (Root)\\Epidermis \& Cortex};

    % Root hairs extending to the right
    \draw[line width=3.5pt, color=agriLeaf!80!black, line cap=round] (2.5, 2.0) .. controls (4.0, 2.2) .. (5.5, 2.0);
    \draw[line width=3.5pt, color=agriLeaf!80!black, line cap=round] (2.5, -1.8) .. controls (4.2, -1.6) .. (6.0, -1.8);
    \node at (4.2, 2.4) [font=\tiny\bfseries\color{agriLeaf!80!black}] {ขนราก (Root Hair)};

    % Rhizosphere Zone boundary (gradient background)
    \draw[color=agriMint, line width=1.5pt, dashed] (6.5, -3.5) -- (6.5, 3.5);
    \node at (4.5, 3.2) [font=\scriptsize\bfseries\color{agriGreen}] {บริเวณไรโซสเฟียร์ (Rhizosphere)};
    \node at (8.5, 3.2) [font=\scriptsize\bfseries\color{slateGrey}] {ดินทั่วไป (Bulk Soil)};

    % Root exudates arrows
    \draw[-{Stealth[scale=1.1]}, line width=1.2pt, color=agriGold] (2.5, 0.8) -- (4.2, 0.8)
        node[midway, above, font=\tiny\color{agriGold!80!black}] {สารหลั่งราก (Sugars, Amino acids)};
    \draw[-{Stealth[scale=1.1]}, line width=1.2pt, color=agriGold] (2.5, -0.6) -- (4.2, -0.6)
        node[midway, below, font=\tiny\color{agriGold!80!black}] {ฟลาโวนอยด์กระตุ้น};

    % PGPR bacteria clusters
    \foreach \x/\y in {2.6/1.4, 2.7/1.7, 2.8/1.2, 3.5/2.1, 4.8/1.9, 2.6/-1.3, 2.7/-2.2, 3.8/-1.7} {
        \draw[fill=agriGreen, draw=white, rounded corners=3pt] (\x,\y) rectangle ++(0.35, 0.18);
    }
    \node at (3.6, 1.4) [font=\tiny\bfseries\color{agriGreen}] {PGPR Biofilm};

    % Signals from bacteria to plant (IAA, Siderophores)
    \draw[-{Stealth[scale=1.1]}, line width=1.2pt, color=agriBlue] (3.8, 1.2) -- (2.5, 0.2)
        node[midway, above right, font=\tiny\bfseries\color{agriBlue}] {ผลิตฮอร์โมน IAA};
    \draw[-{Stealth[scale=1.1]}, line width=1.2pt, color=agriRed] (4.0, -1.2) -- (2.5, -0.2)
        node[midway, below right, font=\tiny\bfseries\color{agriRed}] {ACC Deaminase ลดเครียด};

    % Mycorrhizal hyphae branching into bulk soil
    \draw[line width=1.4pt, color=agriEarth, dashed] (2.5, -3.0) .. controls (4.5, -2.8) and (6.5, -3.2) .. (9.5, -2.5);
    \draw[line width=1.2pt, color=agriEarth, dashed] (6.0, -2.9) .. controls (7.5, -1.8) .. (9.2, -1.5);
    \node at (8.0, -2.2) [font=\tiny\bfseries\color{agriEarth}] {เส้นใยไมคอร์ไรซา (AMF)};
    \node at (8.8, -3.0) [font=\tiny\color{agriBlue}] {ดูดซับ P \& น้ำ};
\end{tikzpicture}
\caption{ภาคตัดขวางรากพืชและปฏิสัมพันธ์เชิงบวกของจุลินทรีย์บริเวณไรโซสเฟียร์ (Rhizosphere Interactions)}
\label{fig:rhizosphere_diagram}
\end{figure}

\section{กลไกการส่งเสริมการเจริญเติบโตของพืช}
\index{กลไกของพีจีพีอาร์ (PGPR Mechanisms)}

PGPR ส่งเสริมการเจริญเติบโตของพืชผ่าน 2 กลุ่มกลไกหลัก:
\begin{enumerate}
    \item \textbf{กลไกทางตรง (Direct Mechanisms):}
    \begin{itemize}
        \item \textbf{การผลิตไฟโตฮอร์โมน (Phytohormone Production)} แบคทีเรียใช้กรดอะมิโนทริปโตเฟน (L-Tryptophan) จากสารหลั่งรากเป็นสารตั้งต้นในการสังเคราะห์ฮอร์โมนออกซิน (Indole-3-acetic acid: IAA) ซึ่งช่วยกระตุ้นการแบ่งเซลล์และการยืดยาวของเซลล์ราก ส่งผลให้พืชแตกรากแขนงและขนรากเพิ่มขึ้น
        \item \textbf{การสร้างเอนไซม์ ACC Deaminase} ในสภาวะเครียด (แล้ง เค็ม หรือน้ำท่วม) พืชจะผลิตฮอร์โมนเอทิลีนในปริมาณสูงเกินไปจนยับยั้งการเติบโต แบคทีเรียที่มีเอนไซม์ 1-aminocyclopropane-1-carboxylate (ACC) deaminase จะเปลี่ยนสารตั้งต้น ACC ให้กลายเป็นแอมโมเนียและแอลฟา-คีโตบิวทีเรต ทำให้ระดับเอทิลีนในพืชลดลงสู่ภาวะสมดุล
        \item \textbf{การสร้างสารไซเดอโรฟอร์ (Siderophores)} เป็นสารประกอบคีเลตน้ำหนักโมเลกุลต่ำที่มีความจำเพาะในการจับกับไอออนเหล็ก ($Fe^{3+}$) สูงมาก ช่วยให้พืชดูดซึมธาตุเหล็กได้ในดินที่มีความเป็นด่าง
    \end{itemize}
    \item \textbf{กลไกทางอ้อม (Indirect Mechanisms):} การแก่งแย่งธาตุอาหารและพื้นที่ผิวราก การหลั่งสารต้านจุลชีพ และการกระตุ้นให้พืชสร้างภูมิคุ้มกันความต้านทานทั่วทั้งต้น (Induced Systemic Resistance: ISR)
\end{enumerate}

\section{ชีววิทยาของเชื้อราอาร์บัสคูลาร์ไมคอร์ไรซา}
\index{ไมคอร์ไรซา (Arbuscular Mycorrhizal Fungi)}

เชื้อราอาร์บัสคูลาร์ไมคอร์ไรซา (Arbuscular Mycorrhizal Fungi: AMF) ในไฟลัม Glomeromycota เป็นเชื้อราที่เจริญอยู่ร่วมกับรากของพืชบกมากกว่า 80\% ของสปีชีส์พืชทั่วโลก เส้นใยราจะแทงผ่านเซลล์ผิวรากเข้าไปยังคอร์เทกซ์และพัฒนาโครงสร้างแตกแขนงคล้ายต้นไม้เล็กๆ เรียกว่า \textbf{อาร์บัสคูล (Arbuscule)} ซึ่งเป็นตำแหน่งที่มีการแลกเปลี่ยนสารอาหารอย่างมีประสิทธิภาพ โดยพืชจะส่งน้ำตาลเฮกโซสและกรดไขมันให้แก่เชื้อรา แลกเปลี่ยนกับฟอสฟอรัส ไนโตรเจน สังกะสี และน้ำที่เส้นใยราดูดซับมาจากดินระยะไกล

\begin{agribox}{การประยุกต์ใช้หัวเชื้อ PGPR และ AMF ในไม้ผลเศรษฐกิจ}
การใส่หัวเชื้อผสมระหว่าง \textit{Bacillus amyloliquefaciens} และเชื้อราไมคอร์ไรซา \textit{Rhizophagus irregularis} ในขั้นตอนการเตรียมหลุมปลูกทุเรียนและมังคุด ช่วยเพิ่มอัตราการรอดตายของต้นกล้าหลังการย้ายปลูกจาก 65\% เป็นมากกว่า 95\% ทั้งยังช่วยให้ระบบรากเจริญเติบโตแผ่กว้างขึ้น ทำให้ต้นพืชทนทานต่อภาวะฝนทิ้งช่วงในฤดูแล้งได้อย่างมีประสิทธิภาพ
\end{agribox}

\section*{คำถามทบทวนท้ายบทที่ 3}
\addcontentsline{toc}{section}{คำถามทบทวนท้ายบทที่ 3}
\begin{enumerate}
    \item สารหลั่งจากรากพืช (Root Exudates) ส่งผลให้เกิดความแตกต่างระหว่างไรโซสเฟียร์กับดินทั่วไปอย่างไร
    \item จงอธิบายบทบาทของเอนไซม์ ACC Deaminase ในการช่วยให้พืชรอดชีวิตภายใต้สภาวะแล้งและดินเค็ม
    \item โครงสร้าง Arbuscule ของเชื้อราไมคอร์ไรซามีความสำคัญอย่างไรในการแลกเปลี่ยนสารอาหารระหว่างพืชและรา
\end{enumerate}

\clearpage